\subsection{Continuum Background for Area--Flow Reduction}
\label{subsec:continuum-constitutive-assumptions}

Only a small amount of continuum background is needed before introducing the
reduced model. At the macroscopic arterial scale, blood-flow models treat
velocity, pressure, density, and stress as continuum fields on a time-dependent
lumen, rather than resolving individual cells. This is the scale at which 3D CFD,
FSI, and 1D area--flow models can be compared as different mathematical tiers;
small-vessel or low-shear regimes require separate rheology caution
\cite{KopplHelmig2023DimensionReducedModeling,GaldiEtAl2008HemodynamicalFlows,PriesEtAl1992BloodViscosityTubeFlow}.

Let $I_T=(0,T_{\mathrm{final}})$ and
$\ITBar=[0,T_{\mathrm{final}}]$. The blood domain at time $t$ is denoted
$\OmgBt$, with space-time cylinder
$\mathcal Q_T=\{(t,\vx):t\in I_T,\ \vx\in\OmgBt\}$. The reference and current
lumen configurations are $\OmgB(0)$ and $\OmgBt$. The notation is summarized in
Appendix~\ref{app:mathematical-notation}; the introduction uses it only to make
clear what is being reduced.

\begin{assumption}[Blood-dynamics kinematic setting]
    \label{ass:blood-dynamics-kinematic-setting}
    Blood motion is represented by a sufficiently regular flow map
    $\La_t:\OmgB(0)\to\OmgBt$ carrying material labels to current positions.
    The Eulerian velocity field $\vu$ satisfies
    \begin{flalign*}
        \quad \pd_t\La_t(\vxi)=\vu(t,\La_t(\vxi)), &&
    \end{flalign*}
    and the accompanying continuum fields are pressure $p$, density $\rho$, and
    Cauchy stress $\vT$ on the current lumen.
\end{assumption}

Mass and momentum balance, together with a stress closure, provide the continuum
source for the lower-dimensional models. For the present report, the essential
point is not the proof of the transport identities; it is that a 1D model stores
cross-sectional quantities and must close the effects of the discarded transverse
velocity and wall response. A section area, volumetric flow, mean axial velocity,
wall law, profile factor, friction law, and rheology descriptor are therefore
part of the reduced model specification rather than incidental implementation
details.

This background also fixes the role of rheology. Large-artery studies often use a
Newtonian approximation, while generalized-Newtonian descriptors can be used when
the effective viscosity is part of the declared closure. In either case, the
comparison in this report depends on the selected closure contract and on the
observation operator, not on an unstated continuum model.
